\clearpage\thispagestyle{empty}
\pdfbookmark[0]{Sources for the review of Zach on mountain attraction}{zachmontagnes-sources}
\begin{center}\Large\bfseries Sources for the review of Zach on mountain attraction\end{center}

This AI translation follows Carl Friedrich Gauss’s review in \emph{Werke}, vol.~VI (Göttingen, 1874), pp.~569--574. The first printing appeared in \emph{Göttingische gelehrte Anzeigen}, no.~210, 31 December 1814, pp.~2089--2100. The tables follow the collected edition’s decimal points and spacing. The form Minet in the density comparison likewise follows the collected edition.

The title page of Franz Xaver von Zach’s book gives Avignon, Seguin the elder, printer and bookseller, 1814. The copy linked below bears the stamp of the Royal Academy of Sciences of Turin and is digitized by the Biblioteca Europea di Informazione e Cultura. Zach describes the hermitage at Notre-Dame des Anges on pp.~343--344 and identifies the three additional sites on p.~399. He discusses the size of the deflection on pp.~356 and 424 and the limits of the longitude observations on pp.~424 and 426--427. The units in the review, including the historical Meile and toise, are retained.

\noindent\textbf{The title cited in the review.} The attraction of mountains and its effects on the plumb lines or levels of astronomical instruments, established and determined by astronomical and geodetic observations made in 1810 at the hermitage of Notre-Dame des Anges on Mount Mimet and at the lighthouse on the island of Planier near Marseille, followed by a geometrical description of the city of Marseille and its territory. By Baron de Zach.

First printing (\href{https://images.p.sub.uni-goettingen.de/iiif/image/gdz:PPN319721507_1814:00002099/full/1650,/0/default.jpg}{2089}, \href{https://images.p.sub.uni-goettingen.de/iiif/image/gdz:PPN319721507_1814:00002100/full/1650,/0/default.jpg}{2090}, \href{https://images.p.sub.uni-goettingen.de/iiif/image/gdz:PPN319721507_1814:00002101/full/1650,/0/default.jpg}{2091}, \href{https://images.p.sub.uni-goettingen.de/iiif/image/gdz:PPN319721507_1814:00002102/full/1650,/0/default.jpg}{2092}, \href{https://images.p.sub.uni-goettingen.de/iiif/image/gdz:PPN319721507_1814:00002103/full/1650,/0/default.jpg}{2093}, \href{https://images.p.sub.uni-goettingen.de/iiif/image/gdz:PPN319721507_1814:00002104/full/1650,/0/default.jpg}{2094}, \href{https://images.p.sub.uni-goettingen.de/iiif/image/gdz:PPN319721507_1814:00002105/full/1650,/0/default.jpg}{2095}, \href{https://images.p.sub.uni-goettingen.de/iiif/image/gdz:PPN319721507_1814:00002106/full/1650,/0/default.jpg}{2096}, \href{https://images.p.sub.uni-goettingen.de/iiif/image/gdz:PPN319721507_1814:00002107/full/1650,/0/default.jpg}{2097}, \href{https://images.p.sub.uni-goettingen.de/iiif/image/gdz:PPN319721507_1814:00002108/full/1650,/0/default.jpg}{2098}, \href{https://images.p.sub.uni-goettingen.de/iiif/image/gdz:PPN319721507_1814:00002109/full/1650,/0/default.jpg}{2099}, \href{https://images.p.sub.uni-goettingen.de/iiif/image/gdz:PPN319721507_1814:00002110/full/1650,/0/default.jpg}{2100}).

\href{https://preserver.beic.it/delivery/DeliveryManagerServlet?dps_pid=IE6956522}{1814 book}, \href{https://preserver.beic.it/delivery/DeliveryManagerServlet?dps_func=download&dps_pid=FL6956565}{Title page}. Selected book pages (\href{https://preserver.beic.it/delivery/DeliveryManagerServlet?dps_func=download&dps_pid=FL6958558}{343}, \href{https://preserver.beic.it/delivery/DeliveryManagerServlet?dps_func=download&dps_pid=FL6958565}{344}, \href{https://preserver.beic.it/delivery/DeliveryManagerServlet?dps_func=download&dps_pid=FL6958645}{356}, \href{https://preserver.beic.it/delivery/DeliveryManagerServlet?dps_func=download&dps_pid=FL6958890}{399}, \href{https://preserver.beic.it/delivery/DeliveryManagerServlet?dps_func=download&dps_pid=FL6959045}{424}, \href{https://preserver.beic.it/delivery/DeliveryManagerServlet?dps_func=download&dps_pid=FL6959057}{426}, \href{https://preserver.beic.it/delivery/DeliveryManagerServlet?dps_func=download&dps_pid=FL6959064}{427}).

\clearpage
